% Összefoglalás magyarul (1500–2500 karakter) – a szerző írja.
A biztonsági műveleti központok elemzőit a behatolásérzékelő rendszerek nagy mennyiségű riasztással terhelik, amelyeknek csak kis része jelez valódi támadást. A túlterhelés riasztási fáradtsághoz vezet, amelynek következtében a valódi támadások is rejtve maradhatnak. A dolgozat célja egy olyan rendszer kidolgozása, amely a hálózati és a gépi behatolásérzékelő riasztásait mérőszámok és a környezetre vonatkozó információ alapján csoportosítja és rangsorolja, együttműködő agentek segítségével.

A szakirodalom áttekintése azt mutatta, hogy a riasztások csoportosítására és rangsorolására javasolt módszerek mindegyike a probléma egy részét oldja meg. A legerősebb, tanítóadatot nem igénylő módszer, a kockázatalapú riasztás-pontozás kizárólag a riasztások folyamából dolgozik, és nem veszi figyelembe az érintett gép fontosságát, sérülékenységeit és megfigyeltségét. A környezetet használó módszerek többnyire szűrnek, így a hiányos nyilvántartás miatt valódi támadást is eltávolíthatnak, a nyelvi modellekre épülő megoldások pedig lassúak és nem megismételhetők.

A választott megoldás a kockázatalapú pontozást három környezeti mérőszámmal egészíti ki, amelyek csak a riasztások sorrendjét módosítják, riasztást nem távolítanak el. A mérőszámokat súlyozott összeggel vonom össze, a súlyokat páros összehasonlítással, a mérés előtt határozom meg. A tervezett rendszer hét agentből áll; a pontszámot és a sorrendet kiszámítható programrészek adják, a helyben futó nyelvi modell csak ellenőrzött magyarázatot ír. A hatást lépésenként bővített modellváltozatokkal, egy nyilvános riasztás-adathalmaz nyolc forgatókönyvén és egy saját laboratóriumi környezetben mérem, a torzítás ellen előzetes rögzítéssel és ellenpróbával védekezve.

A következő félév feladata a rendszer megvalósítása, a mérések elvégzése és az eredmények értékelése.